\chapter{Ambiente di sviluppo}
\label{env}

Questo capitolo descrive l’ambiente di sviluppo e simulazione adottato per il lavoro di tesi, con l’obiettivo di chiarire il ruolo dei principali strumenti software e il motivo della loro combinazione. L’attenzione è posta sull’architettura concettuale della toolchain (simulazione a eventi discreti, simulatore di rete, simulatore di traffico, estensioni per platooning e modellazione cellulare 5G), più che sui dettagli operativi di installazione e compilazione, che saranno trattati nel capitolo successivo.

Nel contesto del platooning cooperativo, la validità degli esperimenti dipende dalla possibilità di rappresentare in modo congiunto fenomeni appartenenti a domini diversi: dinamica veicolare, controllo longitudinale, mobilità stradale e comunicazioni wireless. L’ambiente adottato in questa tesi è stato quindi selezionato per supportare una simulazione accoppiata e modulare, in linea con la letteratura consolidata su simulazione veicolare e cooperative driving \cite{varga2008overview,sommer2011bidirectionally,segata2014plexe,nardini2020simu5g}.

In particolare, il capitolo introduce: i principi della simulazione a eventi discreti per reti veicolari; l’architettura e il workflow sperimentale di OMNeT++; l’integrazione Veins--SUMO tramite TraCI; il ruolo di PLEXE nella simulazione del platooning; il contributo di Simu5G nella modellazione della rete cellulare 5G; e infine le scelte metodologiche adottate per favorire la riproducibilità degli esperimenti.
